
\begin{thebibliography}{*****}




%\bibitem[AL1]{AL-horseshoe} A. Avila and M. Lyubich. 
%  The full renormalization horseshoe for unimodal maps of higher
 % degree: exponential contraction along hybrid classes.
%Publ. Math. IHES 114 (2011), 171-223. 

\bibitem[AL1]{AL-AMS} A. Avila and M. Lyubich. Hausdorff dimension and conformal measures of Feigenbaum Julia sets. J. of the AMS, 21 (2008), 305--383.

\bibitem[AL2]{AL-posmeas} A. Avila and M. Lyubich. 
Lebesgue measure of Feigenbaum Julia sets. Annals Math., v. 176 (2022), 1--88.

\bibitem[ALM]{ALM} A. Avila, M. Lyubich, and W. de Melo. Regular or stochastic dynamics in real analytic families of unimodal maps. Invent. math. v 154 (2003), 451--550. 


\bibitem[BC]{BC} X. Buff and A. Cheritat. Quadratic Julia sets of
  positive area. Annals Math., v. 176 (2012), 673--746.

\bibitem[BD]{BD} B. Branner and A. Douady. Surgery on complex polynomials. Holomorphic dynamics, {M}exico, 1986.



\bibitem[BR]{BR}  A. M. Benini and L. Rempe-Gillen. A landing theorem for entire functions with bounded post-singular sets. GAFA, v 30 (2020), 1465–-1530.


\bibitem[BL]{BL} A. M. Benini and M. Lyubich. Repelling periodic points and landing of rays for post-singularly bounded exponential maps, Annales de l’Institut Fourier 64 (2014), no. 4, 1493--1520.



 \bibitem[Ch]{Che} D. Cheraghi. Topology of irrationally indifferent attractors. arXiv:1706.02678.

 \bibitem[CS]{CS} D. Cheraghi and M. Shishikura. Satellite renormalization of quadratic polynomials. arXiv:1509.07843. 


\bibitem[DG] {dg} R. L. Devaney and  L. R. Goldberg. Uniformization of attracting basins for exponential maps, Duke Math. J. 55 (1987), no. 2, 253--266.



\bibitem[DH1]{DH:Orsay}  A. Douady and J. H. Hubbard. 
  \'{E}tude dynamique des polyn\^{o}mes complexes.
 Publication Mathematiques d'Orsay,
 84-02 and 85-04.


\bibitem[DH2]{DH}
 A. Douady and J. H. Hubbard. On the dynamics of polynomial-like maps.
   Ann. Sc. \'{E}c. Norm. Sup., v. 18 (1985), 287 -- 343.


\bibitem[DK]{DK} R. Devaney and M. Krych. Dynamics of exp(z). Ergodic Theory Dynam. Systems 4 (1984), no. 1, 35--52.


\bibitem[DLS]{DLS} D. Dudko, M. Lyubich, and N. Selinger. Pacman renormalization and self-similarity of the Mandelbrot set near Siegel parameters.  Journal of the AMS, 33 (2020), 653--733.



\bibitem[DS]{DS} A. Dudko and S. Sutherland.  On the Lebesgue measure of the Feigenbaum Julia set. Invent. math, v 221 (2020), 167–-202. 
 





\bibitem[EL]{EL} A. Eremenko and M. Lyubich. Dynamical properties of some classes of entire functions, Ann. Inst. Fourier (Grenoble) 42 (1992), no. 4, 989--1020.

\bibitem[Ep]{Ep}  A. Epstein. Towers of Finite Type Complex Analytic Maps. Ph.D. thesis, 1993.




\bibitem[E1]{E1} H. Epstein, Fixed points of composition operators II. Nonlinearity, vol. 2, (1989), 305-310.


\bibitem[E2]{E2} H. Epstein, Fixed points of the period-doubling operator, Lecture notes, Lausanne.
\bibitem[Er]{E}  A. Eremenko. On the iteration of entire functions. Dynamical systems and ergodic theory (Warsaw, 1986), Banach Center Publ., vol. 23, PWN, Warsaw, 1989, pp. 339--345.


\bibitem[GY]{GY} D. Gaidashev and M. Yampolsky. Renormalization of almost commuting pairs. arXiv:1604.00719. 


\bibitem[H]{Hu} J. H. Hubbard. Local connectivity of Julia sets and bifurcation loci: three theorems of J.-C. Yoccoz. Topological Methods in Modern Mathematics, Proceedings of a Symposium in Honor of John Milnor's Sixtieth Birthday, Publish or Perish, Houston, TX, 1993. 




\bibitem[IS]{IS} H. Inou and M. Shishikura.
  The renormalization for parabolic fixed points and their
  perturbations. Manuscript 2008. 
  
  
\bibitem[K]{K} J. Kahn. A priori bounds for some infinitely renormalizable quadratics: \RN{1}. Bounded primitive combinatorics. Preprint Stony Brook, \# 5 (2006).


\bibitem[J]{J}  Y. Jiang. Infinitely renormalizable quadratic Julia sets. Trans. of the AMS, v 352 (2000), 5077--5091. 


 \bibitem[KL1]{KL1}  J. Kahn and M. Lyubich. A priori bounds for some infinitely renormalizable quadratics: \RN{2}. Decorations. Annals Sci. Ecole Norm. Sup., v. 41 (2008), 57--84.
 
  \bibitem[KL2]{KL2}  J. Kahn and M. Lyubich. A priori bounds for some infinitely renormalizable quadratics, \RN{3}. Molecules. In: ``Complex Dynamics: Families and Friends''.  Proceeding of the conference dedicated to Hubbard's 60th birthday 
(ed.: D. Schleicher ). Peters, A K, Limited, 2009. 



\bibitem[L1]{L:TypBeh} M. Lyubich. Typical behaviour of trajectories of a rational mapping of the sphere. Dokl. Akad. Nauk SSSR, v. 268 (1982), 29 - 32.


\bibitem[L2]{L:acta} M. Lyubich. Dynamics of quadratic polynomials, \RN{1}-\RN{2}. 
Acta Math., v. 178 (1997), 185-297.

\bibitem[L3]{L:FCT} M. Lyubich. Feigenbaum-Coullet-Tresser Universality and Milnor's Hairiness Conjecture.
   Ann. Math., v. 149 (1999), 319 - 420.

\bibitem[L4]{L-book} M. Lyubich. Conformal Geometry and Dynamics of
  Quadratic Polynomials. Book in preparation, \href{http://www.math.stonybrook.edu/~mlyubich/book.pdf}{www.math.sunysb.edu/$\sim$mlyubich/book.pdf}.



\bibitem[LP]{LP} L. Lomonaco and C. L. Petersen. On quasi-conformal (in-) compatibility of satellite copies of the Mandelbrot set: \RN{1}. Invent. math. v 210 (2017), 615--644.

\bibitem[McM1]{McM2} C. McMullen. Renormalization and three manifolds
  which fiber over the circle.   Princeton University Press, 1996. 

\bibitem[McM2]{McM3} C. McMullen. Self-similarity of Siegel disks and Hausdorff dimension of Julia sets.
  Acta Math., v. 180 (1998), 247--292.  


\bibitem[Mi]{Mi} J. Milnor. John Local connectivity of Julia sets: expository lectures. The Mandelbrot set, theme and variations, 67--116, London Math. Soc. Lecture Note Ser., 274, Cambridge Univ. Press, Cambridge, 2000.

%\bibitem[MN]{MN} N.S. Manton, M. Nauenberg. Universal scaling behaviour for iterated maps in the complex plane.
%Commun. Math. Phys. 89(1983), 555-570.


%\bibitem[MP]{MP} R.S. MacKay, I.C. Persival. Universal small-scale structure near the boundary of Siegel disks of
%arbitrary rotation numer. Physica 26D(1987), 193-202.


\bibitem[Na]{Na}  S. B. Nadler, Jr. Continuum Theory. An introduction.
Monographs and Textbooks in Pure and Applied Mathematics, 158, Marcel Dekker, New York, 1992.

\bibitem[PT]{PT} J. Palis and F. Takens. Hyperbolicity and sensitive chaotic dynamics at homoclinic bifurcations. Cambridge Studies in Advanced Mathematics, Vol 35. Fractal dimensions and infinitely many attractors. 1993.



\bibitem[Pe]{Pe} C. Petersen. Local connectivity of some Julia sets containing a circle with an irrational rotation, Acta. Math.,177, 1996, 163-224.



\bibitem[Re]{Re} L. Rempe. Rigidity of escaping dynamics for transcendental entire functions.  Acta Math. 203 (2009), no 2, 235 --267.



\bibitem[RRRS]{RRRS} G. Rottenfusser, J. R\"uckert, L. Rempe, and D. Schleicher. Dynamic rays
of bounded-type entire functions, Ann. of Math. (2) 173 (2011), no. 1, 77--125.


\bibitem[S]{Sh:HD} M. Shishikura. The Hausdorff dimension of the boundary of the Mandelbrot set and Julia sets. Annals Math. (2) 147 (1998), 225--267.


\bibitem[SZ]{SZ} D. Schleicher and J. Zimmer. Escaping points of exponential maps, J. London
Math. Soc. (2) 67 (2003), no. 2, 380--400.

\bibitem[SY]{SY} M. Shishikura and F. Yang. The high type quadratic Siegel disks are Jordan domains. arXiv:1608.04106.



%\bibitem[S]{S} D. Sullivan. Bounds, quadratic differentials and renormalization conjetures.
 % AMS Centennial Publ. 2: Mathematics into Twenty-first century, 1992. 



%\bibitem[ST]{ST} D. Sullivan and W. Thurston. Extending holomorphic motions. Acta Math. 157, 243--257 (1986).





 
\bibitem[Y]{Yarrington} B. Yarrington. Local connectivity and Lebesgue measure of polynomial Julia sets. Thesis,
Stony Brook, 1995.







  




\end{thebibliography}
